\ssscp{Consonants}{03-02-01-01}
Several of the languages of Flores have voiced coronal segments
which are typologically uncommon within Wallacea and/or have
a variety of realisations that require discussion before the
relationships of these languages is discussed.

Firstly, Lio has a voiced coronal affricate/fricative transcribed
by \citet{Elias2018} as ‹ʤ›
and ‹j› and described (p.
32f) as having “[{\ldots}] a number of allophones in free variation
with each other, with no clear conditioning.
The fricative portion of the affricate may be post-alveolar or
palatal, as in [dʒ] or [dʑ],
or it may be closer to a true alveolar, as in [dz].
Furthermore, the initial stop may be entirely gone,
leaving only a voiced fricative such as [z,
ʑ]”\it{.} \citet[6]{Arndt1933} transcribes the same segment ‹dz›
and states: “In some places the \it{dz} sounds as it is written,
like German \it{ds} [dz].
It is widened in some regions towards \it{dsch} [ʤ].
However, nowhere have I heard it so wide that it would correspond
to the German \it{dsch} [ʤ].”\footnote{
		“\it{dz} lautet sich in
		manchen Strichen so wie es geshcrieben wird,
		wie deutsches \it{ds}, verbreitert sich in anderen Gegenden nach
		\it{dsch} hin; doch habe ich es nirgends so breit gehört,
		dass dem deutschen \it{dsch} entspräche.”}
We follow \citet{Elias2018}
by transcribing this segment  ‹ʤ›.

Similarly, Ngadha has a voiced coronal affricate/fricative which
is given by \citet[114]{Djawanai1983} as a voiced alveolar fricative [z].
\citet[9]{Arndt1961} transcribes this segment as ‹dz›
and states: “\it{dz} varies in pronunciation in different regions
between German \it{ds} and \it{dsch}.
It is also heard as \it{z}”.\footnote{
		%Nearly all words in \citet{Arndt1961}
		%with \it{z} have variants with \it{ʣ} and or \it{r}.
		%This variation appears to be due to dialectal variation.
		%Arndt states regarding ‹dz›: “\it{dz} varies in pronunciation
		%in different regions between German \it{ds} and \it{dsch}.
		%It is also heard as \it{z}”.
		%(“\it{dz} schwankt in der Aussprache in verschiedenen Gebieten
		%zwischen deutschen \it{ds} und \it{dsch}; auch hört \it{z} man.”)
		%\citet{Djawanai1983,Djawanai1995} gives only \it{z},
		%described as a voiced alveolar fricative.}
		“\it{dz} schwankt in der
		Aussprache in verschiedenen Gebieten zwischen deutschen \it{ds}
		und \it{dsch}; auch hört \it{z} man.”}
For some words \citet{Arndt1961}
only gives forms with ‹dz›,
for some he gives forms only with ‹z›
and for some he gives variant forms with ‹dz› {\tl} ‹z›.
We follow the transcription used by Arndt in our presentation of Ngadha data.

Rembong also has a voiced alveolar fricative/affricate.
\citet[vi]{Verheijen1977} states rather ambiguously regarding this
segment: “For the sibilant laying between (dz)
and (z) I preferred to take the /z/.” We understand this to mean
that the realisation of the Rembong segment varies between an
alveolar affricate or fricative,
similar to the situation described for Lio and Ngadha.
We follow \citet{Verheijen1977} in transcribing this segment ‹z›.

Several languages of Flores have an alveolar approximant /ɹ/
which contrasts with a trill /r/.
Languages in which this segment is known to occur include: Rongga,
Ende, Watumite Nga{\Q}o, and East Nage.
In none of these languages is this segment in contrast with /z/.
\citet{ArkaEtAl2007} transcribe Rongga /ɹ/ as ‹zh›, while \citet{AokiNakagawa1993}
transcribe Ende /ɹ/ as ‹rh›.
In Rongga, Watumite Nga{\Q}o,
and East Nage this segment develops from PMP *d/*z/*j (see \srf{03-02-02}),
while Ende has *l > \it{ɹ}.

The two varieties of Nga{\Q}o for which data are available have
a segment described by \citet[78]{Elias2018} as a “[{\ldots}]
very unusual phoneme which I transcribe /ɮ/
and which may be described as an optionally pre-stopped voiced palatal lateral fricative.
Before /i/, this phoneme tends towards a realization as an approximant
[l], but in other environments the pre-stopped portion may be dominant
and it approaches [d].” We transcribe this segment ‹d͡ʎ̝›;
that is as a pre-stopped palatal lateral fricative.

Finally, Komodo has contrastive implosives with both plain
voiced /b/ vs. implosive /ɓ/, and /d/ vs. /ɗ/.
However, these are not distinguished in the word list provided
by \citet{Verheijen1987} who writes: “In the end I gave up the
attempt to include the opposition between these phonemes in the orthography.
I myself had difficulty distinguishing the implosives from the
other [plosives] and did not have enough good consultants at my disposal.
It seemed better to me to include no information rather than
unreliable information.” \citep[27]{Verheijen1982}.\footnote{
		“Uiteindelijk heb ik echter de poging om deze twee foneemparen in de orthographie
		op te nemen moeten opgeven.
		Zèlf kon ik de geïmplodeerde klanken maar moeilijk van de andere
		onderscheiden en goede informanten had ik ook niet voldoende ter beschikking.
		Het lijkt mij beter géén dan niet-betrouwbare informatie te verschaffen.”}